\section{Related Work}
\label{sec:related}

\textbf{Model Merging and Task Editing.} Model merging has emerged as a standard practice for consolidating specialized neural models into a single unified architecture without the prohibitive cost of multi-task retraining. Early approaches like FedAvg \cite{mcmahan2017communication} perform simple coordinate-wise parameter averaging. More recently, task arithmetic \cite{ilharco2022editing, ortiz2020task} demonstrated that fine-tuned weight vectors (task vectors) can be linearly added or subtracted to edit model behaviors. To address interference and sign conflicts, advanced techniques have been introduced, such as TIES-Merging \cite{yadav2023ties}, RegMean \cite{jin2022regmean}, Fisher Merging \cite{matena2021merging}, DARE \cite{yu2023dare}, ZipIt! \cite{stoica2023zipit}, and Git Re-Basin \cite{ainsworth2022git}. While highly effective for static deployment, these methods are fundamentally static and cannot adapt to dynamically changing test-time environments.

\textbf{Dynamic Test-Time Merging and Routing.} To handle dynamic query distributions, recent work has explored ensembling specialized parameters on-the-fly. Mixture of Experts (MoE) \cite{shazeer2017outrageously} and Switch Transformers \cite{fedus2021switch} route tokens or samples through specialized sub-networks using trained parametric gating layers. However, MoEs require massive training datasets and complex routing regularization to avoid expert collapse. In contrast, Parameter-Free Subspace Routing (PFSR) \cite{pfsr} introduces a non-parametric approach that derives routing coefficients by projecting features onto low-dimensional subspaces using pre-trained classifier heads. This eliminates the need for a trainable router or calibration split, aligning with minimalist design principles. Unfortunately, PFSR and other parameter-space ensembling methods suffer from severe degradation under mixed-task batch environments, a failure mode we term \emph{heterogeneity collapse}.

\textbf{Systems-centric Stream Robustness.} To address heterogeneity collapse, researchers have developed systems-centric wrappers. Micro-Batch Homogenization (MBH) \cite{mbh} is the state-of-the-art framework in this space. MBH intercepts the query stream and dynamically groups similar tasks using temporal buffering and sorting algorithms to form homogeneous micro-batches. While MBH rescues parameter-space routers, it introduces considerable complexity, stateful serving states, and queuing latency. SABLE replaces this entire systems-level pipeline with a pure, stateless mathematical formulation in activation space, eliminating systems dependencies.

\textbf{Parameter-Efficient Fine-Tuning.} Parameter-Efficient Fine-Tuning (PEFT) methods adapt large pre-trained models using a fraction of the parameters. Standard techniques include adapters \cite{houlsby2019parameter, pfeiffer2020adapterhub, tay2022transporter, hou2022multi}, prompt tuning \cite{lester2021power, dou2022multi}, prefix tuning \cite{li2021prefix}, and BitFit \cite{zaken2022bitfit}. Low-Rank Adaptation (LoRA) \cite{hu2021lora} is a widely adopted PEFT method that parameterizes weight updates as a product of two low-rank matrices. LoRA is typically merged directly into the base weights for zero-overhead inference. In SABLE, we leverage the low-rank properties of LoRA ($r=8$) not just for parameter efficiency during training, but to enable extremely efficient, parallel activation-space ensembling during inference. This elegant design keeps the compute footprint negligible while natively overcoming heterogeneity collapse.

\textbf{Dynamic PEFT Ensembling and Multi-Tenant Serving.} Relatedly, several frameworks have proposed ensembling PEFT adapters. LoraHub~\cite{huang2023lorahub} performs gradient-free optimization to find a static linear combination of LoRA adapters for a given target task, but it requires a target-task calibration split and is static at test-time. MoE-Adapters~\cite{moeadapters} and other routing-based PEFT ensembling methods train a parametric routing network to dynamically weight and blend adapter features. However, these parametric routers require a heavy multi-task training phase, and are highly sensitive to OOD queries. For multi-tenant serving, frameworks like Punica~\cite{punica} and S-LoRA~\cite{slora} compile and run multiple independent adapters in parallel. While these frameworks optimize CUDA kernel dispatch, they serve adapters independently per-sample, and do not provide any ensembling capability. SABLE uniquely bridges this gap: it provides a completely non-parametric, calibration-free dynamic ensembling layer that combines the mathematical soundness of task-vector model merging with the high-speed execution of parallel multi-tenant PEFT serving engines, natively resolving streaming heterogeneity without training or systems-level scheduling buffers.
